\documentclass[10pt,twocolumn]{article}

\usepackage[T1]{fontenc}
\usepackage{lmodern}
\usepackage[margin=0.62in,columnsep=0.24in]{geometry}
\usepackage{amsmath,amssymb,mathtools,bm}
\usepackage{xcolor}
\usepackage{microtype}
\usepackage[colorlinks=true,linkcolor=blue!45!black,
  citecolor=blue!45!black,urlcolor=blue!45!black]{hyperref}

\setlength{\parindent}{0.9em}
\setlength{\parskip}{0pt}
\setlength{\abovedisplayskip}{5pt}
\setlength{\belowdisplayskip}{5pt}
\setlength{\abovedisplayshortskip}{3pt}
\setlength{\belowdisplayshortskip}{3pt}
\setlength{\tabcolsep}{3.2pt}
\raggedbottom

\newcommand{\ket}[1]{\lvert #1\rangle}
\newcommand{\bra}[1]{\langle #1\rvert}
\newcommand{\diag}{\operatorname{diag}}
\newcommand{\argmax}{\operatorname*{arg\,max}}
\newcommand{\chapterfront}[3]{%
  \begingroup
  \centering
  {\small\scshape Integrated Methods Volume\par}
  \vspace{0.45em}
  {\Large\bfseries #1\par}
  \vspace{0.25em}
  {\large #2\par}
  \vspace{0.65em}
  \begin{minipage}{0.94\columnwidth}
    \small\itshape #3
  \end{minipage}\par
  \vspace{0.75em}
  \hrule
  \vspace{0.75em}
  \endgroup
}

\title{SR-SNAKE, JR-SNAKE, and Formal-Manifold SNAKE\\
\large Supported Geometry, Trusted Selection, and Manifold Propagation}
\author{}
\date{14 July 2026}

\begin{document}
\maketitle
\vspace{-1.5em}

\begin{abstract}
This volume presents the three adaptive variational routes separately.
SR-SNAKE uses a three-stage singleton controller: scalar trusted screening in
Phases I and II, followed by a Schur-relaxed, supported-metric Phase-III
response. JR-SNAKE replaces the terminal singleton contest by a joint
active-plus-batch response. Formal-Manifold SNAKE uses JR-SNAKE for structural
acquisition and then propagates the enlarged ansatz through transported tangent
frames, curvature, and adaptive trust. The separately tested FM--SR
composition uses SR-SNAKE as its acquisition law. The route's geometric novelty
appears at its operative phase. Metric support, whitening, rank feasibility,
and trust-region safeguards define the physical coordinate chart used by the
numerical models.
\end{abstract}

\section*{Reading contract and evidence boundary}

The objective is a standalone synthesis of three controllers and their evidence
boundaries. Common geometry is defined once. Each route then receives its own
selection object and transition law. The strong SR curves in the six-regime
comparison identify the
results-backed SR geometry: the controller begins at Phase I, exposes singleton
children, admits one record per round, whitens the supported Phase-III metric,
and carries adaptive trust with beam continuation.
JR supplies the joint-response acquisition law used inside the primary FM
composition. FM--SR remains a visible test composition: SR supplies its
acquisition selector, and FM supplies its formal-manifold propagation law.

The historical result routes also contained ranking and structural policy
decorations that are intentionally removed here. Consequently, ``results-
backed'' means that the energy model, metric convention, candidate exposure,
and propagation mechanism are inherited from those records. Exact historical
adaptive paths additionally depend on shortlist cutoffs, numerical ties, and
the preserved optimizer trajectory.

\newpage
\chapterfront
{Chapter I}
{Common mathematical language}
{This chapter defines the many-body state, candidate records, projective
tangent metric, supported whitening, trust-region quadratic, and optimizer
coordinate charts used later. These shared numerical primitives support
route-specific selectors and route-specific mutable controller state.}

\section*{Variational state and candidate records}

Let the truncated many-body Hilbert space be
\(\mathcal H\simeq\mathbb C^N\). At outer adaptive iteration \(k\), an ordered ansatz
\(U_k(\bm\theta)\) defines
\begin{equation}
\ket{\psi_k}=U_k(\bm\theta_k)\ket{\psi_{\rm ref}},\qquad
E_k=\bra{\psi_k}H\ket{\psi_k}.
\label{eq:state-energy}
\end{equation}
A candidate-position record is \(r=(A,p)\): the Hermitian generator \(A\)
is proposed at circuit position \(p\). If \(U_{>p}\) is the suffix after
that position, then
\begin{align}
\widetilde A_r&=U_{>p}AU_{>p}^{\dagger},\nonumber\\
\tau_r&=-i\widetilde A_r\ket{\psi_k},\qquad
\Pi_k=I-\ket{\psi_k}\bra{\psi_k},\nonumber\\
g_r&=\left|\bra{\psi_k}i[\widetilde A_r,H]\ket{\psi_k}\right|.
\label{eq:candidate-record}
\end{align}
The position is part of the record because the suffix changes the physical
tangent even when the abstract generator label is unchanged.

\section*{Pullback metric, support, and whitening}

For a real parameter displacement
\(\delta\bm\theta\in\mathbb R^n\), let
\begin{equation}
T=[\Pi\partial_1\psi,\ldots,\Pi\partial_n\psi]
:\mathbb R^n\longrightarrow T_{[\psi]}\mathcal M
\label{eq:tangent-map}
\end{equation}
be the horizontal tangent map. Its columns are physical infinitesimal changes
of the quantum ray. The entries of \(\delta\bm\theta\) are coordinate
coefficients that combine those columns. Expanding their physical length
gives
\begin{align}
\|T\delta\bm\theta\|_{\mathbb R}^2
&=\Re\left\langle
\sum_a\delta\theta_a\Pi\partial_a\psi
\middle|
\sum_b\delta\theta_b\Pi\partial_b\psi
\right\rangle\nonumber\\
&=\sum_{a,b}\delta\theta_a
\Re\langle\Pi\partial_a\psi|\Pi\partial_b\psi\rangle
\delta\theta_b\nonumber\\
&=\delta\bm\theta^{\mathsf T}G\delta\bm\theta,
\qquad
G=\Re(T^\dagger T)=V\Lambda V^{\mathsf T}.
\label{eq:raw-gram}
\end{align}
Thus \(G\) is the coordinate representation of physical tangent length. A
null vector of \(G\) is a parameter mixture with zero horizontal state motion.
Numerical damping preserves that physical nullity.

Let \(\lambda_{\max}^+\) be the largest positive raw eigenvalue and retain
\begin{equation}
\mathcal I_+=\{a:\lambda_a>\tau_G\lambda_{\max}^+\}.
\label{eq:support-test}
\end{equation}
If \(V_+\) contains the corresponding eigenvectors and \(\Lambda_+\) their
positive eigenvalues, then
\begin{equation}
GV_+=V_+\Lambda_+,\qquad
V_+^{\mathsf T}V_+=I,\qquad
\Lambda_+\succ0.
\label{eq:supported-spectrum}
\end{equation}
The support retains the certified positive raw modes before any ridge is added.
A materially negative raw eigenvalue is a failed geometry certificate. A valid
stabilized solve begins from a compatible positive-semidefinite raw support.

The exact raw orthonormalizer and its associated physical tangent frame are
\begin{equation}
O_G=V_+\Lambda_+^{-1/2},
\qquad
\mathcal E=TO_G=TV_+\Lambda_+^{-1/2}.
\label{eq:exact-frame}
\end{equation}
Their orthonormality follows directly from the Gram identity:
\begin{align}
\mathcal E^*\mathcal E
&=\Lambda_+^{-1/2}V_+^{\mathsf T}T^*TV_+
\Lambda_+^{-1/2}\nonumber\\
&=\Lambda_+^{-1/2}V_+^{\mathsf T}GV_+
\Lambda_+^{-1/2}=I.
\label{eq:exact-frame-proof}
\end{align}
Define the analysis map
\(L_G=\Lambda_+^{1/2}V_+^{\mathsf T}\). The exact whitened coefficient
\(u\in\mathbb R^{|\mathcal I_+|}\) and its minimum-norm supported lift obey
\begin{equation}
u=L_G\delta\bm\theta,\qquad
\delta\bm\theta_+=O_Gu,\qquad
T\delta\bm\theta_+=\mathcal Eu.
\label{eq:analysis-synthesis}
\end{equation}
Consequently,
\begin{equation}
\delta\bm\theta_+^{\mathsf T}G\delta\bm\theta_+
=u^{\mathsf T}u
=\|\mathcal Eu\|_{\mathbb R}^2.
\label{eq:exact-whitened-length}
\end{equation}
Analysis converts raw parameter coefficients into coefficients along
orthonormal physical modes. Synthesis returns one supported raw-coordinate
representative. The physical frame maps the same coefficients to the tangent
arrow itself. Analysis, synthesis, and frame evaluation therefore connect the
three spaces in Eq.~\eqref{eq:analysis-synthesis}.

\section*{What is whitened, and for how long}

Whitening preserves the quantum state and energy landscape and represents the
retained raw parameter mixtures as coefficients of orthonormal
\emph{physical tangent} arrows at one specified ray. At that ray, the same
infinitesimal motion has three compatible descriptions:
\begin{equation}
u=L_G\delta\bm\theta,
\qquad
\delta\bm\theta_+=O_Gu,
\qquad
T\delta\bm\theta_+=\mathcal Eu.
\label{eq:three-whitening-descriptions}
\end{equation}
The first expression analyzes a raw parameter mixture, the second chooses its
minimum-norm supported parameter representative, and the third displays the
physical tangent arrow. The coefficient \(u\) is the coordinate list of that
tangent arrow in the orthonormal frame \(\mathcal E\).

Every object in Eq.~\eqref{eq:three-whitening-descriptions} is local to the
base ray and coordinate registry used to construct \(T\). During one model
solve that source-point chart is held fixed, so the differential, curvature,
trust constraint, and lifted step all refer to the same tangent plane. Moving
the ray changes the tangent columns; adding coordinates changes their registry.
Either event requires a new raw Gram matrix, a new support test, and a new
whitening frame. SR reconstructs outer selection geometry for the next
structural round. FM connects consecutive inner frames so that a learned
curvature model can be represented at the new endpoint.

The numerical solve may condition the retained inverse response with a metric
ridge \(\mu_G\geq0\). Define
\begin{align}
C_G&=\Lambda_+^{1/2}(\Lambda_++\mu_GI)^{-1/2},\nonumber\\
W_G&=V_+(\Lambda_++\mu_GI)^{-1/2}=O_GC_G.
\label{eq:three-metric-maps}
\end{align}
If \(x\) is the stabilized solver coordinate, then the complete coordinate
chain is
\begin{equation}
\delta\bm\theta_+=W_Gx=O_Gu,
\qquad
u=C_Gx.
\label{eq:stabilized-coordinate-chain}
\end{equation}
The two induced lengths are therefore
\begin{align}
O_G^{\mathsf T}GO_G&=I,\nonumber\\
W_G^{\mathsf T}GW_G&=C_G^{\mathsf T}C_G,\nonumber\\
\|T\delta\bm\theta_+\|_{\mathbb R}^2
&=x^{\mathsf T}C_G^{\mathsf T}C_Gx.
\label{eq:metric-identities}
\end{align}
The exact coefficient \(u\) measures raw Fubini--Study tangent length. The
stabilized coefficient \(x\) is the coordinate in which the numerical trust
solve is performed. They coincide when \(\mu_G=0\). For \(\mu_G>0\), the
bridge \(C_G\) records the distinction explicitly, preventing a ridge from
being mistaken for a change in physical distance.

Whitening must also respect a change of parameter chart. Let
\(\bm\theta=S\bm\phi\) with nonsingular Jacobian \(S\), so that
\(\delta\bm\theta=S\delta\bm\phi\). The chain rule gives
\begin{equation}
T_\phi=T_\theta S,\qquad
G_\phi=S^{\mathsf T}G_\theta S.
\label{eq:chart-metric-transform}
\end{equation}
The same tangent and its length satisfy
\begin{equation}
T_\phi\delta\bm\phi=T_\theta\delta\bm\theta,\qquad
\delta\bm\phi^{\mathsf T}G_\phi\delta\bm\phi
=\delta\bm\theta^{\mathsf T}G_\theta\delta\bm\theta.
\label{eq:chart-length-invariance}
\end{equation}
For a downhill covector \(d\) and ordinary Hessian \(Q\), consistent chart
representations obey
\begin{equation}
d_\phi=S^{\mathsf T}d_\theta,\qquad
Q_\phi=S^{\mathsf T}Q_\theta S.
\label{eq:chart-model-transform}
\end{equation}
These transformations preserve both
\(d^{\mathsf T}\delta\bm\theta\) and
\(\delta\bm\theta^{\mathsf T}Q\delta\bm\theta\). A rescaling of angle units
therefore changes coordinate arrays, tangent columns, Gram entries,
differential components, and curvature components together. The physical
update and its predicted energy response are invariant under this rescaling.

\section*{Trusted quadratic response}

Let \(d\in\mathbb R^n\) be a downhill differential and
\(Q=Q^{\mathsf T}\in\mathbb R^{n\times n}\) an ordinary energy Hessian.
For the supported raw displacement \(\delta\bm\theta=W_Gx\), define
\begin{align}
\widehat{\Delta E}(x)
&=d_x^{\mathsf T}x-\tfrac12x^{\mathsf T}Q_xx,\nonumber\\
d_x&=W_G^{\mathsf T}d,\qquad
Q_x=W_G^{\mathsf T}QW_G.
\label{eq:supported-quadratic}
\end{align}
The route solves this model on \(\|x\|_2\leq\rho\). The lifted raw physical
length \(\|C_Gx\|_2\), and when required the exact endpoint
Fubini--Study distance, remain separate certificates. A curvature shift
\(\sigma_H\geq0\) may be required on the retained support. With
\(A_x=Q_x+\sigma_HI\), the free supported step is
\begin{equation}
x_0=A_x^{-1}d_x.
\label{eq:free-step}
\end{equation}
If \(\|x_0\|_2\leq\rho\), the completed trust test returns a zero multiplier.
For \(\|x_0\|_2>\rho\), choose \(\lambda_T>0\) so that
\begin{equation}
x_T=(A_x+\lambda_TI)^{-1}d_x,\qquad \|x_T\|_2=\rho.
\label{eq:boundary-step}
\end{equation}
The radius is route-local controller state. Chapter II defines one branch-local
SR radius, frozen across every
metric-bearing phase in a controller round. Chapter IV defines an FM radius
inside the post-growth manifold optimizer. Their respective update laws
preserve the meanings of displacement agreement and energy-model agreement.

For a one-dimensional supported direction with metric length \(F>0\),
downhill differential \(g\geq0\), curvature \(h\), and radius \(\rho\),
the trusted decrease is obtained directly from
\begin{equation}
\max_{0\leq\alpha\leq\rho/\sqrt F}
\left(g\alpha-\tfrac12[h]_+\alpha^2\right),
\qquad [h]_+=\max(h,0).
\label{eq:scalar-trust-response}
\end{equation}
The three SR phases below display their own substitutions directly inside this
one-dimensional problem.

\section*{Logical generators and the Powell base chart}

A logical generator can expand into several runtime Pauli rotations and retain
one intended variational angle. Let logical block \(\ell\) contain
runtime indices \(I_\ell\), with \(m_\ell=|I_\ell|\). Write
\(\bm\phi\in\mathbb R^d\) for logical angles and
\(\bm\vartheta\in\mathbb R^m\) for runtime coordinates. Define the block
average \(P\in\mathbb R^{d\times m}\) and uniform lift
\(R\in\mathbb R^{m\times d}\) by
\begin{equation}
P_{\ell j}=\begin{cases}m_\ell^{-1},&j\in I_\ell,\\0,&\text{otherwise},\end{cases}
\qquad
R_{j\ell}=\begin{cases}1,&j\in I_\ell,\\0,&\text{otherwise}.
\end{cases}
\label{eq:logical-runtime-maps}
\end{equation}
Then \(PR=I_d\), and \(RP\) projects a runtime block to its uniform
representative. State preparation in the shared-angle convention depends on
\begin{equation}
\bm\phi=P\bm\vartheta,\qquad
\ket{\psi(\bm\vartheta)}=\ket{\psi(RP\bm\vartheta)}.
\label{eq:runtime-projection}
\end{equation}
Hence the expanded runtime chart contains flat alias directions
\(\ker P\). The physical shared-angle ansatz family remains unchanged.

The SR refit begins from this expanded-runtime/projected-logical base chart.
Powell's underlying parameter array has one entry for each executable Pauli
factor. Before every circuit evaluation, \(P\) replaces the entries within
each logical block by their mean and \(R\) copies that mean back to the
runtime factors. The exact state therefore depends on
\(RP\bm\vartheta\). Chapter II derives the supported Fubini--Study map that is
placed around this base chart before Powell receives its optimization
variables.

\section*{Controller modules}

It is useful to name four maps as distinct stages:
\begin{equation}
\begin{aligned}
\text{expose}&\longrightarrow\text{select}
\longrightarrow\text{zero growth}\\
&\longrightarrow\text{enlarged-manifold optimization}.
\end{aligned}
\label{eq:four-modules}
\end{equation}
SR-SNAKE specifies the first two maps and ordinary zero-angle growth. FM-SNAKE
specifies a richer final map and can be paired with more than one selector.
This modular reading resolves the informal phrase ``FM with SR inside'': SR
chooses the structural enlargement; formal-manifold propagation optimizes the
enlarged ansatz.

\clearpage
\chapterfront
{Chapter II}
{Results-backed SR-SNAKE}
{This chapter uses the three-phase singleton route behind the strong SR curves,
retains the route's phase-specific geometric novelty, sets the cost factor to
unity, and follows the controller through its Phase-III singleton decision and
ensuing transition. The Phase-III singleton score supplies terminal authority.}

\section*{Route identity}

At one branch and outer iteration \(k\), the controller is
\begin{equation}
\mathcal R_k^{(1)}\longrightarrow
\mathcal R_k^{(2)}\longrightarrow
\mathcal R_k^{(3)}\longrightarrow r_k^*.
\label{eq:sr-route}
\end{equation}
The controller begins at Phase I. Each \(r=(A,p)\) remains a singleton
generator-position record. Macro parents may expose symmetry-legal singleton
Pauli children. Phases II and III select from singleton records. Each beam
branch evaluates its own state-dependent record population.

\section*{Phase I: local metric trust}

Let
\begin{equation}
F_r=\Re\langle\Pi_k\tau_r,\Pi_k\tau_r\rangle.
\label{eq:sr-F-raw}
\end{equation}
The first screen replaces an unknown directional Hessian by the metric-scaled
proxy \(\lambda_FF_r\), giving
\begin{equation}
\Delta E_1(r)=
\max_{0\leq\alpha\leq\rho_k^{\mathrm S}/\sqrt{F_r}}
\left(g_r\alpha-\tfrac12\lambda_FF_r\alpha^2\right).
\label{eq:sr-phase1}
\end{equation}
The first screen combines derivative magnitude with Fubini--Study length. A
direction with a large derivative and large physical tangent length can receive
a smaller trusted response.

\section*{Phase II: measured directional curvature}

Phase II measures the candidate directional energy curvature \(h_r\). With a
small scalar metric stabilization \(\epsilon_F\geq0\), let
\begin{equation}
F_{2,r}=\max(F_r+\epsilon_F,F_{\min}).
\label{eq:sr-phase2-metric}
\end{equation}
The Phase-II trusted energy response is
\begin{equation}
\Delta E_2(r)=
\max_{0\leq\alpha\leq\rho_k^{\mathrm S}/\sqrt{F_{2,r}}}
\left(g_r\alpha-\tfrac12[h_r]_+\alpha^2\right).
\label{eq:sr-phase2}
\end{equation}
Let \(T_{A,k}\) be the active tangent map,
\(G_{A,k}=\Re(T_{A,k}^\dagger T_{A,k})\), and
\(q_{2,r}=\Re(T_{A,k}^\dagger\Pi_k\tau_r)\). The route-specific Phase-II
novelty is the residual tangent fraction
\begin{equation}
\mathcal N_2(r)=
1-\frac{q_{2,r}^{\mathsf T}G_{A,k}^{+}q_{2,r}}{F_r},
\qquad
\mathcal V_2(r)=\Delta E_2(r)\mathcal N_2(r).
\label{eq:sr-phase2-novelty}
\end{equation}
The retained population \(\mathcal R_k^{(3)}\) is determined by
\(\mathcal V_2\), hard feasibility, and deterministic tie rules. This is the
singleton novelty used by SR. Batch novelty receives its own definition in the
joint-response route.

\section*{Phase III: Schur-relaxed curvature and metric}

Let \(W_r\) be the active coordinate window used for the final singleton
model. Denote its Gram and ordinary energy-Hessian blocks by
\(Q_r\) and \(H_r\), respectively. Let \(q_r\) contain active--candidate
tangent overlaps, \(b_r\) the active--candidate energy-Hessian column,
\(F_r\) the raw candidate metric, and \(h_r\) its raw curvature. The active
relaxation vector \(u_r\) solves
\begin{equation}
(H_r+\mu_HI)u_r=b_r,
\label{eq:sr-window-relaxation}
\end{equation}
on a certified stable branch. The relaxed scalar curvature and the metric
length of the same relaxed displacement are
\begin{align}
h_r^*&=h_r-b_r^{\mathsf T}u_r,\nonumber\\
F_r^*&=F_r-2q_r^{\mathsf T}u_r+u_r^{\mathsf T}Q_ru_r,\nonumber\\
q_r^*&=q_r-Q_ru_r.
\label{eq:sr-relaxed-pair}
\end{align}
These quantities must be read together. The Hessian solve says how the active
coordinates compensate the proposed direction. The Gram expression then
measures the physical length of that compensated direction. A candidate is
infeasible if the raw supported geometry fails or if \(F_r^*\) collapses
below its support floor.

The supported-metric solver implements this reduction in a whitened chart and
checks the lifted Fubini--Study displacement. Its scalar Schur reading is
\begin{equation}
\Delta E_3(r)=
\max_{0\leq\alpha\leq\rho_k^{\mathrm S}/\sqrt{F_r^*}}
\left(g_r\alpha-\tfrac12[h_r^*]_+\alpha^2\right).
\label{eq:sr-phase3}
\end{equation}
The compensated tangent is projected against the same active span through
\(q_r^*\). Its surviving fraction and the terminal SR score are
\begin{equation}
\mathcal N_3(r)=
1-\frac{(q_r^*)^{\mathsf T}Q_r^+q_r^*}{F_r^*},
\qquad
\mathcal V_3(r)=\Delta E_3(r)\mathcal N_3(r).
\label{eq:sr-phase3-novelty}
\end{equation}
The terminal decision is therefore
\begin{equation}
\boxed{
r_k^*\in\argmax_{r\in\mathcal R_k^{(3)}}\mathcal V_3(r).}
\label{eq:sr-decision}
\end{equation}
This is the SR Phase-III model: Schur-relaxed singleton gain followed by the
route's explicit singleton novelty multiplier. The decision law uses
\(\mathcal V_3(r)\) explicitly. A full active-plus-singleton conditional
difference such as \(D_{\rm S}(r)\) defines another mathematical construction.

\section*{Location of matrix whitening in SR}

Phases I and II already carry trust through the scalar metric lengths in
Eqs.~\eqref{eq:sr-phase1} and \eqref{eq:sr-phase2}. Phase III is the first
stage that solves a coupled active-coordinate model and therefore needs a
matrix support factorization and whitening map. Phases I and II use scalar
trust. Phase III uses supported matrix trust in the coupled final stage.

Coordinate rescaling safety transforms every coupled object through the same
map: Gram, differential, Hessian, trust displacement, and lifted seed. The
identities in Eq.~\eqref{eq:metric-identities} provide the numerical
certificates for this consistent chart transformation.

\section*{Lifetime of the outer SR whitening chart}

At outer round \(k\), every Phase-III coupled model is constructed at the
entrance ray \([\psi_k]\), before the winning record is inserted. Its supported
chart simultaneously represents the active compensation, candidate motion,
differential, curvature, and trust displacement. The chart remains frozen
during the solve and comparison of the surviving records. This is
why a whitening transformation must be applied to the complete coupled model
including \(G_k\), the differential, the curvature, and the displacement.

The resulting Phase-III step predicts a tangent displacement and supplies the
lifted proposal associated with the selected record. After zero-coordinate
insertion, the subsequent nonlinear refit constructs
a full-ansatz supported Fubini--Study frame around the Powell base chart from
Chapter I. Powell iterates in that accepted-refit frame. This is a new frame at
the accepted zero-growth state, independently rebuilt from the full accepted
ansatz after the Phase-III active-plus-candidate decision. The exact post-refit
ray is used afterward to
measure the realized Fubini--Study displacement and update
\(\rho_{k+1}^{\mathrm S}\).

The next outer round reconstructs its state-dependent selection Gram and its
accepted-refit Gram at their respective anchors. Each round begins with a fresh
coordinate frame and fresh optimizer state. The accepted-refit orthonormalizer
remains fixed for one Powell invocation and is rebuilt after the next admission
changes the ansatz. Chapter IV uses repeated endpoint frames differently
because FM carries curvature information between accepted inner moves.

\section*{The fixed supported-FS Powell chart}

\subsection*{The accepted zero-growth origin}

Suppose round \(k\) has admitted a generator at zero angle. The coordinate
registry has enlarged at the accepted pre-refit quantum ray. Let
\(\bm\phi_k^{(0)}\in\mathbb R^{d_k}\) contain all accepted logical
angles and let \(\bm\vartheta_k^{(0)}\in\mathbb R^{m_k}\) be their expanded
runtime representation. With the block maps from
Eq.~\eqref{eq:logical-runtime-maps},
\begin{equation}
\bm\vartheta_k^{(0)}=R_k\bm\phi_k^{(0)},
\qquad
\ket{\psi_k^{(0)}}
=U_k(\bm\vartheta_k^{(0)})\ket{\psi_{\rm ref}}.
\label{eq:sr-refit-origin}
\end{equation}
The refit includes the full accepted ansatz. Thus \(d_k\) counts every
accepted logical angle and extends beyond the smaller active window used by the
selector. The newly admitted coordinate begins at zero angle and contributes
its tangent column.

An expanded runtime displacement
\(\delta\bm\vartheta\in\mathbb R^{m_k}\) is converted to a logical
displacement and then returned to the uniform runtime representation:
\begin{equation}
\delta\bm\phi=P_k\delta\bm\vartheta,
\qquad
\widehat{\bm\vartheta}_k(\delta\bm\vartheta)
=R_k\bigl(\bm\phi_k^{(0)}+P_k\delta\bm\vartheta\bigr).
\label{eq:sr-refit-runtime-map}
\end{equation}
Powell's base array is expanded because it contains one position per
executable factor. Circuit preparation uses the projected logical block means
\(P_k\delta\bm\vartheta\) as the shared logical angles. The
hat on \(\widehat{\bm\vartheta}_k\) denotes the runtime array actually sent to
the circuit executor.

\subsection*{The metric seen by the base chart}

Let \(T_k^{\phi}:\mathbb R^{d_k}\rightarrow
T_{[\psi_k^{(0)}]}\mathcal M_k\) be the horizontal tangent map with respect to
the logical angles. The tangent map associated with the expanded base array is
the composition
\begin{equation}
T_k^{\vartheta}=T_k^{\phi}P_k.
\label{eq:sr-refit-base-tangent}
\end{equation}
Consequently, if \(G_k^{\phi}=\Re[(T_k^{\phi})^*T_k^{\phi}]\), then the Gram
matrix in the base coordinates is
\begin{align}
G_k^{\vartheta}
&=\Re[(T_k^{\vartheta})^*T_k^{\vartheta}]\nonumber\\
&=P_k^{\mathsf T}G_k^{\phi}P_k.
\label{eq:sr-refit-base-gram}
\end{align}
This equation is the pullback of the logical Fubini--Study metric through the
block-average map. It also exposes the alias directions. For any
\(a\in\ker P_k\),
\begin{equation}
P_ka=0
\quad\Longrightarrow\quad
T_k^{\vartheta}a=0
\quad\Longrightarrow\quad
G_k^{\vartheta}a=0.
\label{eq:sr-refit-alias-null}
\end{equation}
Such an array changes the numbers stored inside a runtime block and preserves
its mean, the executed circuit, and the quantum ray. The supported
factorization removes these zero-motion combinations before Powell begins.

\subsection*{Whitening the supported runtime directions}

Diagonalize the raw base Gram matrix and retain its positive supported
modes:
\begin{equation}
G_k^{\vartheta}V_{k,+}=V_{k,+}\Lambda_{k,+},
\qquad
V_{k,+}^{\mathsf T}V_{k,+}=I,
\qquad
\Lambda_{k,+}>0.
\label{eq:sr-refit-supported-spectrum}
\end{equation}
If \(r_k\) is the retained rank, define the synthesis map
\begin{equation}
O_k=V_{k,+}\Lambda_{k,+}^{-1/2}
:\mathbb R^{r_k}\longrightarrow\mathbb R^{m_k}.
\label{eq:sr-refit-orthonormalizer}
\end{equation}
The variable supplied to Powell is
\(z\in\mathbb R^{r_k}\). Its complete lift is
\begin{align}
\delta\bm\vartheta&=O_kz,
&\delta\bm\phi&=P_kO_kz,\nonumber\\
\widehat{\bm\vartheta}_k(z)
&=R_k\bigl(\bm\phi_k^{(0)}+P_kO_kz\bigr).
\label{eq:sr-refit-whitened-lift}
\end{align}
The defining certificate is
\begin{align}
O_k^{\mathsf T}G_k^{\vartheta}O_k&=I,\nonumber\\
\left\|T_k^{\vartheta}O_kz\right\|_{\mathbb R}^2
&=z^{\mathsf T}z.
\label{eq:sr-refit-whitened-length}
\end{align}
Thus every unit coordinate axis received by Powell represents a unit-length
physical tangent direction at the refit origin, and distinct axes represent
orthogonal physical tangents. This is the precise sense in which the inner
optimizer is geometrically preconditioned. Whitening changes the numerical
coordinates used by Powell and preserves the energy landscape and the
existing ranking factors.

\subsection*{The objective actually evaluated by Powell}

The complete coordinate chain is
\begin{equation}
z\xmapsto{\ O_k\ }\delta\bm\vartheta
\xmapsto{\ P_k\ }\delta\bm\phi
\xmapsto{\ R_k\ }\widehat{\bm\vartheta}_k(z)
\xmapsto{\ U_k\ }\ket{\psi_k(z)}
\xmapsto{\ H\ }E_k(z).
\label{eq:sr-refit-coordinate-chain}
\end{equation}
In full,
\begin{align}
\ket{\psi_k(z)}
&=U_k\!\left(
R_k\bigl(\bm\phi_k^{(0)}+P_kO_kz\bigr)
\right)\ket{\psi_{\rm ref}},\nonumber\\
E_k(z)&=\bra{\psi_k(z)}H\ket{\psi_k(z)},
\qquad z^{(0)}=0.
\label{eq:sr-refit-powell-objective}
\end{align}
Powell proposes values of \(z\) and receives the corresponding energy. It
manipulates \(z\) exclusively. The classical maps \(O_k\), \(P_k\), and
\(R_k\) convert each trial into the logical and runtime arrays used to prepare
the quantum state.

The starting value \(z^{(0)}=0\) is important. It represents the inherited
accepted state exactly, so Powell's first direction set is centered at the
state whose Gram matrix defined the whitening. Near that origin, equal
Euclidean displacements of \(z\) correspond to equal first-order
Fubini--Study tangent lengths. The energy evaluations then determine which of
those comparably scaled directions produces useful nonlinear descent.

\subsection*{One fixed frame for one Powell invocation}

The Gram matrix \(G_k^{\vartheta}\) and orthonormalizer \(O_k\) are constructed
once at \(\ket{\psi_k^{(0)}}\). Every Powell line search and every trial point
within that invocation reuses the same map in
Eq.~\eqref{eq:sr-refit-whitened-lift}. Trial evaluations therefore require the
energy \(E_k(z)\) and reuse the single anchor Gram construction.

This fixed-frame rule is both the computational boundary and the geometric
approximation. The identity in Eq.~\eqref{eq:sr-refit-whitened-length} is exact
for tangent motion at the origin. At a distant nonlinear trial point, the
local tangent metric may evolve. Powell continues to use the original
coordinate map. The chart acts as a local preconditioner fixed throughout the
finite-budget search.

When the next generator is admitted, the accepted ansatz and its tangent map
change. SR then constructs
\begin{equation}
G_{k+1}^{\vartheta}
\longrightarrow
V_{k+1,+},\Lambda_{k+1,+}
\longrightarrow
O_{k+1}
\label{eq:sr-refit-next-frame}
\end{equation}
at the next zero-growth origin. One full supported-FS factorization is thereby
associated with one accepted-refit Powell invocation. Every energy trial inside
that invocation reuses it.

\subsection*{Why this chart helps Powell}

Before whitening, the physical unit-length surface in the base array is the
ellipsoid
\begin{equation}
\delta\bm\vartheta^{\mathsf T}
G_k^{\vartheta}\delta\bm\vartheta=1.
\label{eq:sr-refit-base-ellipsoid}
\end{equation}
Its long and short axes mean that equal numerical changes of two runtime
coordinates can produce very different state motions. Alias directions extend
flatly because they produce zero state motion. After the supported synthesis
\(\delta\bm\vartheta=O_kz\), the same local physical surface is
\begin{equation}
z^{\mathsf T}z=1.
\label{eq:sr-refit-whitened-sphere}
\end{equation}
Powell's line searches now begin with directions normalized by physical
tangent length. Removal of the raw scale disparity and exact alias directions
allows the finite direction set to spend its evaluations on energy descent.
Energy curvature remains free to differ from direction to direction. The chart
equalizes local state-motion scale, and the objective determines energetic
stiffness. The scientific ansatz and objective are invariant under this
coordinate transformation. The numerical
path through a finite Powell budget can improve substantially because the
optimizer receives a better-conditioned representation of that same accepted
manifold.

\section*{SR trust-radius update}

SR carries one branch-local radius for the complete controller round. If
\(\rho_k^{\mathrm S}\) is frozen at the round entrance, then every
metric-bearing stage receives the same value:
\begin{equation}
\rho_{1,k}^{\mathrm S}
=\rho_{2,k}^{\mathrm S}
=\rho_{3,k}^{\mathrm S}
=\rho_k^{\mathrm S}.
\label{eq:sr-one-radius}
\end{equation}
The subscripts on the left identify three consumption locations for one
branch-local controller state. The update occurs after the winning structural
proposal has undergone its nonlinear refit and before any pruning mutation.

Let \(\widehat z_k\) be the applied Phase-III model step in the
pre-admission coordinate chart, with pre-admission Gram matrix \(G_k\). Its
predicted physical displacement is
\begin{equation}
\widehat\delta_k
=\sqrt{\max(0,\widehat z_k^{\mathsf T}G_k\widehat z_k)}.
\label{eq:sr-predicted-displacement}
\end{equation}
For the scalar Phase-III reading,
\(\widehat\delta_k=|\alpha_k^*|\sqrt{F_{r_k^*}^*}\). Let
\(\ket{\psi_{k+1}^{+}}\) denote the accepted post-refit, pre-prune state. The
realized displacement is
\begin{equation}
\delta_k
=\arccos\!\left(
\left|\langle\psi_k|\psi_{k+1}^{+}\rangle\right|
\right),
\qquad
\chi_k^{\mathrm S}
=\frac{\delta_k}{\widehat\delta_k+\varepsilon_\delta}.
\label{eq:sr-displacement-ratio}
\end{equation}
The overlap makes \(\delta_k\) invariant to global phase. The ratio measures
realized physical motion relative to predicted physical motion.

Let \(c_k=1\) when the previous trust constraint was binding, the exact
post-refit energy decrease exceeds its numerical acceptance tolerance, and any
enabled direction-consistency guard passes. The unbounded displacement-
calibrated update is
\begin{align}
0\leq\chi_k^{\mathrm S}<1
&\Longrightarrow
\rho_{k+1}^{\mathrm S}
=\max\!\left(\rho_\epsilon,
\rho_k^{\mathrm S}\sqrt{\chi_k^{\mathrm S}}\right),
\label{eq:sr-trust-contract}\\
\chi_k^{\mathrm S}>1,\quad c_k=1
&\Longrightarrow
\rho_{k+1}^{\mathrm S}
=\rho_k^{\mathrm S}\sqrt{\chi_k^{\mathrm S}},
\label{eq:sr-trust-expand}\\
\chi_k^{\mathrm S}\geq1,\quad c_k=0
&\Longrightarrow
\rho_{k+1}^{\mathrm S}=\rho_k^{\mathrm S}.
\label{eq:sr-trust-hold}
\end{align}
Here \(\rho_\epsilon\) supplies the positive floating-point representability
floor. The SR radius follows the unbounded multiplicative update in
Eqs.~\eqref{eq:sr-trust-contract}--\eqref{eq:sr-trust-hold}. An invalid or
nonfinite displacement transaction also
holds \(\rho_k^{\mathrm S}\), thereby reserving contraction evidence for valid
physical-displacement measurements. Expansion receives credit from a binding
prior solve whose active radius limited the model step.

The canonical profile starts from \(\rho_0^{\mathrm S}=0.25\) and uses
\(\varepsilon_\delta=10^{-12}\). Its central selection object remains the
relaxed singleton geometry \((h_r^*,F_r^*,q_r^*)\), followed by
\(\mathcal N_3(r)\). The trust update manages the next round's physical scale;
the route's singleton novelty remains \(\mathcal N_3(r)\), and the completed
round retains its selected candidate.

\clearpage
\chapterfront
{Chapter III}
{JR-SNAKE}
{JR-SNAKE uses the Phase-I and Phase-II candidate funnel and a jointly modeled
candidate batch as its terminal selection object. Its terminal authority is the
conditional active-plus-batch response, with route identity defined by that
joint response.}

\section*{Joint-response route identity}

Let \(B\subseteq\mathcal R_k^{(3)}\) be a feasible batch assembled from the
Phase-II survivors, and write
\(z_B=(\delta\bm\theta,\bm\alpha_B)\) for the inherited active-coordinate
response and proposed batch amplitudes. The complete active-plus-batch model
uses its full Gram matrix \(G_B\), ordinary Hessian \(Q_B\), and downhill
differential \(d_B\). Supported whitening supplies \(z_B=W_{G,B}x\).

JR compares the best trusted response in the enlarged coordinate space with
the best response already available from active-coordinate relaxation alone:
\begin{align}
D_{\rm J}(B)={}&
\max_{\substack{z_B=W_{G,B}x\\\|x\|_2\leq\rho_k^{\mathrm J}}}
\left(d_B^{\mathsf T}z_B-\tfrac12z_B^{\mathsf T}Q_Bz_B\right)
\nonumber\\
&-
\max_{\substack{z_B=(\delta\bm\theta,\bm0)\\
z_B=W_{G,B}x,\ \|x\|_2\leq\rho_k^{\mathrm J}}}
\left(d_B^{\mathsf T}z_B-\tfrac12z_B^{\mathsf T}Q_Bz_B\right).
\label{eq:jr-conditional}
\end{align}
The same trust radius and supported coordinate chart appear in both maxima.
The second maximum isolates descent available from active-coordinate
relaxation with \(\bm\alpha_B=\bm0\). Their difference assigns the batch its
conditional contribution. The terminal JR decision is
\begin{equation}
\boxed{
B_k^{\rm J}\in\argmax_{B\in\mathfrak F_k}D_{\rm J}(B),}
\label{eq:jr-decision}
\end{equation}
where \(\mathfrak F_k\) is the route's feasible batch frontier.

The JR acquisition radius \(\rho_k^{\mathrm J}\) is selector-side state. Its
displacement-calibrated update has the same form as
Eqs.~\eqref{eq:sr-predicted-displacement}--\eqref{eq:sr-trust-hold}, with
\(\widehat z_k\) replaced by the selected batch's complete joint step. In an
FM--JR composition this acquisition radius remains distinct from the FM inner
radius \(\rho_j^{\mathrm F}\) defined in Chapter IV.

\section*{Where novelty enters JR}

The records reaching \(\mathfrak F_k\) have already passed the route's
Phase-II singleton novelty, whose geometric form is shown in
Eq.~\eqref{eq:sr-phase2-novelty}. The terminal JR contest then uses
Eq.~\eqref{eq:jr-conditional}; the conditional joint response supplies its
terminal score.

\section*{JR growth and refit}

After Eq.~\eqref{eq:jr-decision}, every member of the selected batch is added
in the same structural transition and the enlarged ansatz is refitted by the
JR continuation policy. The batch frontier, joint response, and batch
admission remain JR state. The next chapter introduces formal-manifold
transport.

\clearpage
\chapterfront
{Chapter IV}
{Formal-Manifold SNAKE}
{The primary FM composition uses the JR acquisition law from Chapter III and
then applies transported formal-manifold propagation after zero-coordinate
growth. The separately tested FM--SR composition uses the SR acquisition law
with the same propagation core.}

\section*{The common propagation core}

Suppose a selector has accepted a singleton or subset \(B_k\). Zero-coordinate
growth produces
\begin{equation}
\bm\theta_k\longmapsto
\bm\theta_{k+1}^{(0)}=(\bm\theta_k,\bm0),
\qquad [\psi_{k+1}^{(0)}]=[\psi_k].
\label{eq:fm-zero-growth}
\end{equation}
The physical ray retains its value. The coordinate registry and available
tangent space can expand. Let \(D_{B_k}\) collect the horizontal tangent
columns created by the accepted coordinates at zero. If \(\mathcal E_k\) is
the orthonormal frame of the old supported tangent space, then
\begin{equation}
T_{k+1}^{(0)}=[T_k,D_{B_k}],\qquad
R_{B_k}=(I-\mathcal E_k\mathcal E_k^*)D_{B_k}.
\label{eq:fm-zero-growth-tangents}
\end{equation}
The columns of \(R_{B_k}\) are the components of the new coordinate tangents
in the orthogonal complement of the old supported frame. Their retained rank
determines whether the enlarged coordinate registry has actually enlarged the
physical tangent quotient. This geometric statement characterizes the new
manifold chart after the selector has completed its decision.

\section*{Outer FM structural frame extension}

Zero-coordinate growth is a structural frame-extension event at an unchanged
ray. The tangent space may have gained independent directions. Resolve the new
residual block through
\begin{equation}
G_{B_k}^{\perp}
=\Re(R_{B_k}^{\dagger}R_{B_k})
=V_{B_k}\Lambda_{B_k}V_{B_k}^{\mathsf T}.
\label{eq:fm-growth-residual-gram}
\end{equation}
After applying the same raw support rule used elsewhere, the retained new
physical frame and enlarged frame are
\begin{equation}
\mathcal E_{B_k}
=R_{B_k}V_{B_k,+}\Lambda_{B_k,+}^{-1/2},
\qquad
\mathcal E_{k+1}^{(0)}
=[\mathcal E_k,\mathcal E_{B_k}].
\label{eq:fm-growth-frame-extension}
\end{equation}
The inherited columns remain the old physical motions. The appended columns
are the supported components of the admitted tangents in the old frame's
orthogonal complement. Residualization establishes their
orthogonality before diagonalization normalizes the new--new block.

In raw parameter coordinates the enlarged Gram matrix has the block form
\begin{equation}
G_{k+1}^{(0)}
=
\begin{pmatrix}
G_{WW}&G_{WB}\\
G_{BW}&G_{BB}
\end{pmatrix}.
\label{eq:fm-growth-gram-blocks}
\end{equation}
Because the ray is unchanged, \(G_{WW}\) can be inherited when the old
coordinate subsequence and parameterization are exactly preserved. The
old--new and new--new blocks contain geometric information about the admitted
directions and must be supplied by the growth receipt or reconstructed at the
same anchor state. A candidate-space Gram from SR is reusable under a complete
active-plus-new block certificate with matching coordinate provenance.

The transition from an \(r\)-dimensional frame to an
\((r+s)\)-dimensional frame uses isometric embedding and frame extension. FM
checks whether the old frame embeds isometrically in the enlarged frame. The
corresponding inverse-curvature prior initializes newly admitted modes with
the inherited-plus-new form
\begin{equation}
B_{k+1}^{(0)}
=
\begin{pmatrix}
B_k&0\\
0&\beta I_s
\end{pmatrix},
\qquad \beta>0.
\label{eq:fm-growth-curvature-extension}
\end{equation}
The zero mixed blocks encode initially unknown old--new energy-curvature
coupling. The first accepted inner secant can begin to learn those couplings.
A failed isometric-embedding or support certificate resets the incompatible
curvature history.

This produces three separate geometric operations in an FM--SR composition.
SR whitens the entrance candidate model to select a structural addition. FM
then rebuilds and, when certified, extends the full active tangent frame at the
unchanged zero-growth ray. Finally, the inner FM loop repeatedly whitens
endpoint active geometry and transports the persistent curvature model during
motion on the fixed enlarged ansatz. This volume follows that fixed-manifold
propagation. A future pruning rule receives its own certified restriction or
reduction.

The outer index \(k\) labels structural growth. The inner index \(j\) begins
after Eq.~\eqref{eq:fm-zero-growth} and labels formal-manifold optimization
steps on the fixed enlarged ansatz. The selector has finished its structural
decision before the inner propagation state is initialized.

\section*{The two iteration clocks}

The outer clock changes which tangent directions are available. It starts from
\([\psi_k]\), evaluates SR or JR entrance geometry, selects a structural
enlargement, and inserts its new coordinates at zero. At that instant the ray
remains \([\psi_k]\). The parameter registry and possible tangent motions have
changed. The outer selector's whitening has completed its decision and is
followed by a newly constructed FM optimizer frame.

The inner clock evolves the ray on that fixed enlarged registry.
At source point \(j\), FM builds \(T_j\), \(G_j\), \(O_j\), and
\(\mathcal E_j\); represents the differential and curvature in that frame;
chooses one trusted step; and holds the source frame fixed through the line
search. The accepted parameter update prepares the endpoint state. FM then
builds \(T_{j+1}\), \(G_{j+1}\), \(O_{j+1}\), and
\(\mathcal E_{j+1}\). Because the curvature model persists across inner
steps, the old model must be transported into the new frame before it can be
corrected by endpoint data.

\section*{Supported endpoint frames and whitening}

At inner point \(j\), let
\begin{align}
T_j&=[\Pi_j\partial_1\psi_j,\ldots,\Pi_j\partial_d\psi_j],\nonumber\\
G_j&=\Re(T_j^\dagger T_j)
=V_j\Lambda_jV_j^{\mathsf T}.
\label{eq:fm-inner-gram}
\end{align}
The raw support test is repeated at every accepted endpoint because the tangent
columns and their linear dependencies depend on the current quantum state. Let
\(V_{j,+}\) and \(\Lambda_{j,+}\) denote the retained positive spectrum. The
exact physical frame is
\begin{equation}
O_j=V_{j,+}\Lambda_{j,+}^{-1/2},
\qquad
\mathcal E_j=T_jO_j.
\label{eq:fm-exact-frame}
\end{equation}
Its orthonormality follows from the endpoint Gram matrix itself:
\begin{equation}
\mathcal E_j^*\mathcal E_j
=O_j^{\mathsf T}G_jO_j=I.
\label{eq:fm-frame-identity}
\end{equation}
Thus \(\mathcal E_j\) contains physical tangent arrows. The map \(O_j\) sends
coefficients in that physical frame back into one supported parameter
representative. Recomputing this frame makes the local model follow the
evolving projective state and supersedes the initial parameter grid.

The regularized bridge and stabilized solver map are
\begin{align}
C_j&=\Lambda_{j,+}^{1/2}
(\Lambda_{j,+}+\mu_GI)^{-1/2},\nonumber\\
W_j&=O_jC_j
=V_{j,+}(\Lambda_{j,+}+\mu_GI)^{-1/2}.
\label{eq:fm-maps}
\end{align}
For a raw parameter displacement \(\delta\bm\theta_j\), exact orthonormal
coefficient \(v\), and stabilized coefficient \(x\),
\begin{equation}
\delta\bm\theta_j=O_jv=W_jx,
\qquad
v=C_jx,
\qquad
T_j\delta\bm\theta_j=\mathcal E_jv.
\label{eq:fm-coordinate-bridge}
\end{equation}
Therefore
\begin{equation}
\|T_j\delta\bm\theta_j\|_{\mathbb R}^2
=\|v\|_2^2
=x^{\mathsf T}C_j^{\mathsf T}C_jx.
\label{eq:fm-coordinate-length}
\end{equation}
The exact coefficient \(v\) measures raw Fubini--Study tangent length. The
stabilized coefficient \(x\) conditions the numerical solve. The bridge
between them is retained as state: a numerical ridge conditions inverse
response, and the raw support fixes the surviving physical directions.

At a later endpoint, \(O_{j+1}\), \(C_{j+1}\), and \(W_{j+1}\) are rebuilt
from \(G_{j+1}\). Even if the retained rank stays fixed, the columns of
\(\mathcal E_{j+1}\) can rotate relative to \(\mathcal E_j\). Consequently,
an array copied unchanged from step \(j\) generally represents a different
physical tangent at step \(j+1\). This is why formal-manifold propagation must
transport both vectors and curvature representations. The transport supplies
geometric continuity across the numerical warm start.

\section*{Quantum data and classical whitening}

Whitening is classical linear algebra applied to state-dependent quantum
information. The classical controller performs multiplication by \(O_j\),
\(C_j\), and \(W_j\). The quantum circuit prepares the parameterized
state. At an anchor state, the controller requires the energy \(E_j\),
parameter differential \(b_j\), and raw tangent Gram matrix
\begin{equation}
(G_j)_{ab}
=\Re\langle\Pi_j\partial_a\psi_j|
\Pi_j\partial_b\psi_j\rangle.
\label{eq:fm-measured-gram-entry}
\end{equation}
On quantum hardware these are obtained from expectation-value or overlap
measurements. An exact statevector calculation obtains the same arrays
classically and thereby consumes the prepared-state information represented by
hardware measurements.

Once those arrays are available, support selection, eigendecomposition,
construction of \(O_j\), the trust solve, line search bookkeeping,
Procrustes singular-value decomposition, curvature transport, inverse-RBFGS or
SR1, and the trust-radius update are classical computations. The resulting
ordinary parameter proposal \(p_j=O_ju_j\) is returned to state preparation.
The quantum device or simulator then prepares trial parameters
\(\bm\theta_j+a_jp_j\) and supplies the data used to judge the endpoint.

The information requirements are different at different gates. Trial energy
is needed to test the line search and energy-model agreement. An accepted
endpoint requires a new differential for the secant, a new Gram matrix for
exact support and rewhitening, and old--new tangent-frame overlaps for
Procrustes transport. The endpoint Gram and cross-frame overlaps require
state-dependent information in addition to the endpoint energy. A symmetric
\(d\times d\) Gram matrix contains
\begin{equation}
\frac{d(d+1)}{2}
\label{eq:fm-gram-entry-count}
\end{equation}
distinct real entries before observable grouping or other measurement
compression. Exact reconstruction after every accepted inner move therefore
has quadratic metric-acquisition scaling in the active dimension.

\section*{Exact reference geometry and a practical anchor policy}

The current FM implementation is an exact-state reference route. Its
authoritative metric is reconstructed from analytic state tangents at every
accepted endpoint. Its qBroyden metric recurrence is shadow telemetry: the
exact endpoint Gram overwrites the prediction, and every metric acquisition
therefore remains exact. Moreover, the current exact backend returns energy,
differential, statevector, and tangent data for every line-search evaluation,
including rejected trials. That behavior is convenient for deterministic
simulation and substantially more information-intensive than a staged hardware
acquisition policy. A hardware realization requires a statistical rank
certificate in place of the route's exact-state rank certificate.

A hardware-facing policy would first screen line-search trials with energy.
When a trial passes that gate, it would acquire the endpoint differential,
metric, and cross-frame overlaps needed to certify and commit the new FM state.
Rejected energy trials could then reuse the source-point \(G_j\), \(b_j\),
frame, and curvature model and vary the trial scale. This staged
policy defines a prospective implementation requirement. The present backend
uses exact-state acquisition.

An anchor-and-recycle metric policy can reduce repeated full-Gram acquisition.
It would measure an authoritative \(G_j\)
at selected accepted inner points and use a labeled predictor between those
anchors. Exact reanchoring would remain mandatory when support or rank is
uncertain, frame alignment fails, metric innovation is too large, conditioning
becomes unsafe, or the trust model disagrees materially with the realized
endpoint. A predicted metric may guide response between anchors. Physical-mode
admission remains governed by the fail-closed support test at each required
anchor.

Consequently, exact per-step rewhitening and recycled inner geometry define two
FM profiles. The implemented profile supplies mathematically exact geometry at
measurement-intensive cost. The scalable design profile requires an
operational predictor, anchor cadence, statistical rank certificates,
invalidation rules, and regression tests to establish its trajectory identity.

\section*{Trusted inverse-curvature propagation}

Let \(b_j=\nabla E(\bm\theta_j)\) be the ordinary parameter differential and
\(\gamma_j=O_j^{\mathsf T}b_j\) its representation in the exact orthonormal
frame. The authoritative FM profile carries a positive inverse-curvature model
\(B_j\) in that frame; write \(A_j=B_j^{-1}\). Its local energy-change model
is
\begin{equation}
m_j(v)=\gamma_j^{\mathsf T}v
+\tfrac12v^{\mathsf T}A_jv.
\label{eq:fm-local-model}
\end{equation}
The free inverse-curvature step solves
\(A_ju_j=-\gamma_j\), hence \(u_j=-B_j\gamma_j\). FM then applies its own
inner trust radius before lifting the step back to parameters:
\begin{align}
u_j&\leftarrow u_j
\min\!\left(1,\frac{\rho_j^{\mathrm F}}{\|u_j\|_2}\right),\nonumber\\
p_j&=O_ju_j.
\label{eq:fm-proposal}
\end{align}
The trust sphere is Euclidean because \(u_j\) is expressed in an orthonormal
physical tangent frame. Directions judged stiff by \(A_j\) receive less
motion. Compliant directions receive more motion. The radius limits how far the
recycled curvature model may extrapolate from the current quantum state.

A line search selects \(a_j\in(0,1]\) and accepts the source-frame displacement
\begin{equation}
\delta_j=a_ju_j,
\qquad
\delta\bm\theta_j
:=\bm\theta_{j+1}-\bm\theta_j
=a_jp_j=O_j\delta_j.
\label{eq:fm-accepted-step}
\end{equation}
Its first-order physical ray displacement at the source is therefore
\begin{equation}
T_j\delta\bm\theta_j=\mathcal E_j\delta_j.
\label{eq:fm-accepted-tangent}
\end{equation}
Here \(u_j\) is the trust-capped proposal in the exact orthonormal frame,
\(p_j\) is its raw parameter lift before line search, \(\delta_j\) is the
accepted orthonormal-frame coefficient after line search, and
\(\delta\bm\theta_j\) is the actual accepted parameter change. The tangent in
Eq.~\eqref{eq:fm-accepted-tangent} is the source-point, first-order description
of the ray motion. The actual endpoint ray is obtained by preparing the circuit
at \(\bm\theta_{j+1}\); its energy follows the nonlinear circuit map and may
differ from the quadratic prediction. After preparing that state, FM rebuilds
its energy, differential, horizontal tangents, raw metric support, and
whitening. These exact endpoint
objects determine whether the proposal may become the next accepted FM state.

\section*{FM trust-radius update}

FM calibrates its inner radius from energy-model agreement. For the accepted
step \(\delta_j\), define predicted and realized decreases by
\begin{align}
\widehat d_j
&=-\left(\gamma_j^{\mathsf T}\delta_j
+\tfrac12\delta_j^{\mathsf T}A_j\delta_j\right),
\nonumber\\
d_j&=E_j-E_{j+1},
\nonumber\\
\chi_j^{\mathrm F}&=\frac{d_j}{\widehat d_j},
\qquad \widehat d_j>0.
\label{eq:fm-agreement-ratio}
\end{align}
This is an actual-to-predicted energy-reduction ratio. It differs from the SR
ratio in Eq.~\eqref{eq:sr-displacement-ratio}, which compares realized and
predicted physical displacement.

The distinction is consequential. The prediction \(\widehat d_j\) evaluates
the source-frame model \(m_j\) along the accepted coefficient \(\delta_j\).
The realization \(d_j\) evaluates the energy after the nonlinear circuit map
has produced the endpoint ray. Their ratio tests the local energy model used by
FM exclusively. The outer SR admission score and SR radius retain their values
from the outer controller.

Let \(c_j^-=1\) when the ratio is unavailable or
\(\chi_j^{\mathrm F}<1/4\). Let \(c_j^+=1\) when
\(\chi_j^{\mathrm F}>3/4\) and the trust-capped, pre-line-search direction
reached at least four fifths of the current trust radius,
\(\|u_j\|_2\geq4\rho_j^{\mathrm F}/5\). The multiplicative update is
\begin{align}
\kappa_j&=
\begin{cases}
\tfrac12,&c_j^-=1,\\
\tfrac32,&c_j^-=0,\ c_j^+=1,\\
1,&c_j^-=c_j^+=0,
\end{cases}
\nonumber\\
\rho_{j+1}^{\mathrm F}
&=\min\!\left(\rho_+^{\mathrm F},
\max\!\left(\rho_-^{\mathrm F},
\kappa_j\rho_j^{\mathrm F}\right)\right).
\label{eq:fm-trust-update}
\end{align}
The current profile uses
\(\rho_0^{\mathrm F}=0.25\),
\(\rho_-^{\mathrm F}=10^{-8}\), and
\(\rho_+^{\mathrm F}=2\). Strong agreement expands the radius when the proposed
step nearly reached the boundary. Boundary-reaching steps supply the evidence
for a larger neighborhood.

A rejected line-search or trust proposal contracts the pending radius through
\begin{equation}
\rho_j^{\mathrm F}\longmapsto
\max\!\left(\rho_-^{\mathrm F},\tfrac12\rho_j^{\mathrm F}\right).
\label{eq:fm-rejection-contraction}
\end{equation}
An accepted endpoint with a supported-rank change, failed frame transport, or
hard metric innovation applies the same contraction and resets the curvature
representation. A failed curvature postcondition applies an
additional half contraction. Valid zero-coordinate growth transports the
committed radius unchanged; incompatible growth contracts it. Rollback
discards the pending proposal and preserves the radius stored in the last
complete accepted FM state.

SR and FM therefore manage different state at different times. SR updates
\(\rho_k^{\mathrm S}\) once between outer controller rounds from displacement
agreement through its unbounded multiplicative law. FM updates
\(\rho_j^{\mathrm F}\) inside the fixed enlarged manifold from energy-model
agreement and clips it to explicit inner-optimizer bounds. An FM--SR
composition may contain both radii: the SR radius governs entrance selection,
and the FM radius governs post-growth propagation. Each controller writes its
own trust state exclusively.

\section*{Endpoint frame transport and curvature recycling}

When the retained rank is unchanged, define the endpoint cross-frame overlap
and its singular-value decomposition by
\begin{align}
K_{j+1,j}&=\mathcal E_{j+1}^*\mathcal E_j
=U_j\Sigma_jV_j^{\mathsf T},
\nonumber\\
Q_j&=U_jV_j^{\mathsf T},
\qquad Q_j^{\mathsf T}Q_j=I.
\label{eq:fm-procrustes}
\end{align}
The orthogonal factor \(Q_j\) is the Procrustes alignment that pairs the two
observed orthonormal frames as closely as their cross overlaps allow. The
corresponding physical transport is
\begin{equation}
\mathcal T_j
=\mathcal E_{j+1}Q_j\mathcal E_j^*.
\label{eq:fm-physical-transport}
\end{equation}
If \(\zeta_j=\mathcal E_jv\) is a tangent at the source, then
\(\mathcal T_j\zeta_j=\mathcal E_{j+1}Q_jv\) is its endpoint counterpart.
A passive coordinate change rewrites a vector within one tangent space;
\(\mathcal T_j\) actively relates two different endpoint tangent spaces.

The same alignment transports direct and inverse curvature by
\begin{equation}
\widetilde A_j=Q_jA_jQ_j^{\mathsf T},
\qquad
\widetilde B_j=Q_jB_jQ_j^{\mathsf T}.
\label{eq:fm-curvature-transport}
\end{equation}
For \(v_+=Q_jv\),
\begin{equation}
v_+^{\mathsf T}\widetilde A_jv_+
=v^{\mathsf T}A_jv.
\label{eq:fm-quadratic-invariance}
\end{equation}
This preserved scalar quadratic response is the reason curvature can be
recycled. Applying \(Q_j\) attaches the old soft and stiff directions to their
corresponding physical endpoint modes.

The order is essential. FM first rebuilds the endpoint whitening frame, then
aligns the old and new physical frames through \(Q_j\), then transports the old
curvature representation, and then uses the accepted endpoint data to
update that transported model. Re-whitening changes the coordinate array used
to describe the learned quadratic response. The secant below supplies one new
observed energy-curvature action after both observations have been placed in
the endpoint frame.

The accepted displacement and differential change must also be compared in
one frame. Their endpoint representations are
\begin{equation}
s_j=Q_j\delta_j,
\qquad
y_j=\gamma_{j+1}-Q_j\gamma_j.
\label{eq:fm-secant}
\end{equation}
The secant pair teaches the action of curvature along the one accepted motion;
the prior model and later secants supply the remaining curvature actions. With
a safeguarded differential change \(\bar y_j\) satisfying the
positive-curvature certificate and
\(\xi_j=(s_j^{\mathsf T}\bar y_j)^{-1}\), inverse Riemannian BFGS gives
\begin{align}
B_{j+1}={}&
(I-\xi_js_j\bar y_j^{\mathsf T})\widetilde B_j
(I-\xi_j\bar y_js_j^{\mathsf T})\nonumber\\
&+\xi_js_js_j^{\mathsf T}.
\label{eq:fm-bfgs-compact}
\end{align}
The damping applies the minimum raw-differential correction required to retain
a positive inverse-curvature model. The accepted endpoint, transported
frame, and guarded secant must all pass their postconditions before this model
replaces the previous one.

In regularized solver coordinates the same physical transport is
\begin{equation}
P_j=C_{j+1}^{-1}Q_jC_j.
\label{eq:fm-regularized-transport}
\end{equation}
Vectors use \(P_j\), covectors use \(P_j^{-\mathsf T}\), direct curvature
uses \(P_j^{-\mathsf T}AP_j^{-1}\), and inverse curvature uses
\(P_jBP_j^{\mathsf T}\). These transformation laws preserve vector--covector
pairings and quadratic response after both the physical frame and the
regularization bridge have changed.

The exact reference route uses inverse-RBFGS for objective curvature and keeps
qBroyden as shadow prediction telemetry. The observed endpoint
Gram remains authoritative and overwrites that prediction. Making the
predictor operational between selected anchors would define the separate
recycled-geometry policy described above. The two compositions below share
the present exact propagation core. Their acquisition selectors act before
Eq.~\eqref{eq:fm-zero-growth}.

\section*{Primary composition: FM--JR}

The primary FM route uses the JR decision
\(B_k^{\rm J}\) from Eq.~\eqref{eq:jr-decision}. Every coordinate in that
selected batch is inserted at zero through Eq.~\eqref{eq:fm-zero-growth}; the
resulting enlarged ansatz is then optimized through
Eqs.~\eqref{eq:fm-inner-gram}--\eqref{eq:fm-regularized-transport}. JR owns
the candidate frontier and the conditional joint-response decision. FM owns
the supported endpoint frames, transported curvature, trust state, and
accepted continuation after growth.

This composition has a precise handoff. JR finishes with a selected structural
batch, its zero-coordinate insertion map, and its acquisition receipts. FM
begins from the resulting unchanged ray and enlarged coordinate registry. It
rebuilds the supported tangent frame before taking an inner step. The ordinary
joint-response Hessian used to choose \(B_k^{\rm J}\) is therefore entrance
selection curvature. FM initializes the transported inverse-curvature state
\(B_j\) from its propagation state.

The two trust radii also remain separate. JR uses
\(\rho_k^{\mathrm J}\) to decide which structural batch to admit. After
the batch is inserted, FM initializes or restores \(\rho_j^{\mathrm F}\) and
updates it by Eq.~\eqref{eq:fm-trust-update}. The FM radius initializes
independently from the JR entrance model, and the completed JR acquisition
radius remains fixed throughout FM line search.

Query closure means that the JR geometry and response used for one decision
refer to the same entrance state, coordinate registry, and parameterization.
It is a consistency requirement attached to the JR acquisition data.

\section*{Test composition: FM--SR}

The tested FM--SR composition replaces the JR batch decision by the SR
singleton decision from Chapter II:
\begin{equation}
r_k^{\mathrm S}\in
\argmax_{r\in\mathcal R_k^{(3)}}\mathcal V_3(r).
\label{eq:fm-sr-selector}
\end{equation}
The selected singleton is inserted at zero and then optimized by the same FM
propagation equations. This test combines SR's staged singleton acquisition
with FM's transport-aware nonlinear optimizer and preserves the SR score and SR
trajectory identity.

The entrance prediction and the FM endpoint have different jobs.
\(\mathcal V_3(r)\) ranks candidate structures at the entrance state. The FM
cycle determines the realized nonlinear endpoint after growth. The SR radius
\(\rho_k^{\mathrm S}\) is frozen across the three entrance phases and is
updated from physical displacement after the accepted SR transition. The
FM radius \(\rho_j^{\mathrm F}\) controls inner propagation after zero growth
and is updated from the local FM energy model. In particular,
\(\chi_j^{\mathrm F}\) uses the predicted decrease from
Eq.~\eqref{eq:fm-local-model} along the actual inner step
\(\delta_j\). The quantity \(\mathcal V_3(r)\) remains the entrance selection
score.

This separation permits a meaningful composition test. A changed FM endpoint
records the effect of the transported post-growth optimizer under fixed SR
candidate ordering, novelty factor, and outer-round trust law. A changed SR
entrance record supplies a new manifold under fixed FM frame-transport and
curvature-learning laws.

\section*{Commit and rollback}

An FM proposal is committed after all endpoint energy, metric support,
whitening, frame transport, curvature, and model-agreement certificates pass.
If any certificate fails, the ansatz point and every mathematical object whose
meaning depends on that point return to their last accepted values. This is one
atomic state transition, expressed directly through the coupled state objects.

Atomicity matters because these objects describe one another. The differential
is meaningful at a particular state; the supported frame is built from that
state's tangent Gram matrix; the curvature array is represented in that frame;
and the trust radius summarizes the reliability of that local model. Retaining
new curvature with an old frame, or a contracted pending radius with a restored
old state, would splice together incompatible geometric facts. A rejected
transaction therefore leaves the complete previously committed FM point
unchanged. A successful transaction advances the state, energy, support,
whitening, transported curvature, and \(\rho_{j+1}^{\mathrm F}\) together.

\section*{The three route boundaries}

SR-SNAKE owns a staged singleton funnel. Its Phase-II and Phase-III novelty
factors act on one candidate-position record, and its terminal decision admits
exactly one record. JR-SNAKE owns a joint candidate frontier. It gives terminal
authority to the conditional active-plus-batch response in
Eq.~\eqref{eq:jr-conditional}; this joint response defines its route identity.

Formal-Manifold SNAKE begins after a structural acquisition law has selected
the enlargement. In the primary composition that acquisition law is JR; in
the tested FM--SR composition it is SR. FM remains a distinct route because it
owns the zero-growth transition, supported endpoint frames, transported
curvature, trust update, and atomic commit or rollback after the ansatz manifold
has changed.

\footnotesize
\begin{thebibliography}{9}

\bibitem{Grimsley2019}
H.~R. Grimsley, S.~E. Economou, E. Barnes, and J.~P. Mayhall,
An adaptive variational algorithm for exact molecular simulations on a
quantum computer, \emph{Nature Communications} \textbf{10}, 3007 (2019).

\bibitem{Stokes2020QNG}
J. Stokes, J. Izaac, N. Killoran, and G. Carleo,
Quantum natural gradient, \emph{Quantum} \textbf{4}, 269 (2020).

\bibitem{Ramoa2025HessianRecycling}
M. Ramoa et al., Reducing measurement costs by recycling the Hessian in
adaptive variational quantum algorithms, \emph{Quantum Science and
Technology} \textbf{10}, 015031 (2025).

\bibitem{Sohail2026GeoADAPT}
M.~A. Sohail and T. Koike-Akino, Geo-ADAPT-VQE: Quantum Information
Metric-Aware Circuit Optimization for Quantum Chemistry,
\emph{arXiv:2603.10325} (2026).

\end{thebibliography}

\end{document}
